\PassOptionsToPackage{expansion=false}{microtype}
\documentclass[11pt,a4paper,sans]{moderncv}
% FONT
\usepackage[T1]{fontenc}
\usepackage{lmodern}
\usepackage{enumitem}
% STYLE
\moderncvstyle{classic}
\moderncvcolor{blue}
% PAGE LAYOUT
\usepackage[scale=0.82]{geometry}
\setlength{\hintscolumnwidth}{2.8cm}
%----------------------------------------------------------------------------------
% PERSONAL DATA
%----------------------------------------------------------------------------------
\name{Camille}{Robichon}
\title{R\&D Formulation Scientist -- Lipid Systems \& Antioxidants}
\address{Université de Montpellier, UMR QualiSud, CIRAD}{Montpellier, France}{}
\phone[mobile]{+33 6 52 80 32 84}
\email{camillerobichon30@gmail.com}
\social[linkedin]{camille-robichon-8a5a46202}
\social[github]{camillerobichon30}
\extrainfo{ORCID: 0009-0007-2489-364X}
\begin{document}
\makecvtitle
%----------------------------------------------------------------------------------
% SUMMARY
%----------------------------------------------------------------------------------
\section{Research Profile}
\cvitem{}{%
Formulation scientist specialised in lipid-based delivery systems, emulsions, liposomes, niosomes and antioxidant strategies. Experienced in translating molecular mechanisms into innovative lipid formulations with improved oxidative stability and bioactive delivery. Skilled
in experimental design, kinetic modelling, and multidisciplinary analytical techniques for food, nutraceutical and cosmetic applications. International research experience in France, the United
Kingdom, and Japan.
}
%----------------------------------------------------------------------------------
% EDUCATION
%----------------------------------------------------------------------------------
\section{Education}
\cventry
{2023--Present}
{PhD in Food Science}
{University of Montpellier}
{Montpellier, France}
{}
{%
{\textbf{Expected completion: October 2026}\\
Thesis: \emph{New Formulation Strategies for Vitamin E to Improve Antioxidant
Efficiency in Food Matrices}.}
\cventry
{2020--2023}
{Chemical Engineering Degree}
{École Nationale Supérieure de Chimie de Lille (ENSCL)}
{Lille, France}
{}
{\textit{Specialisation: Formulation Chemistry.}
}
\cventry
{Sep 2022--Mar 2023}
{Exchange Semester}
{Doshisha University}
{Kyoto, Japan}
{}
{Project: Transdermal delivery of curcumin using an oil-in-water microemulsion.}
%----------------------------------------------------------------------------------
% PROFESSIONAL & RESEARCH EXPERIENCES
%----------------------------------------------------------------------------------
\section{Research Experience}
\cventry
{Nov 2023--Present}
{PhD Candidate in Food Science}
{CIRAD, UMR QualiSud}
{Montpellier, France}
{}
{%
\emph{New Formulation Strategies for Vitamin E to Improve Antioxidant Efficiency in Food Matrices.}
\begin{itemize}[leftmargin=*,nosep]
\item Designed lipid formulations and screened antioxidant combinations to improve vitamin E stability and oxidative resistance in oil-in-water emulsions.
\item Developed and compared Weibull, Gompertz, Logistic and Exponential models to characterise lipid oxidation kinetics and quantify antioxidant efficiency and interaction effects.
\item Developed and optimised lipid model systems and nanoemulsions using microfluidisation to investigate oxidation mechanisms.
\item Formulated and characterised liposomes and niosomes using DLS and HPLC and evaluated their effects on vitamin E stability and bioaccessibility.
\item Monitored lipid oxidation kinetics using NMR, HPLC, GC--MS, and UV--Vis spectroscopy.
\item Investigated the distribution of bioactive compounds between emulsion droplets and surfactant micelles using Taylor dispersion analysis.
\item Evaluated vitamin E bioaccessibility using a static \emph{in vitro} gastrointestinal digestion model.
\item Research outcomes include two first-author publications and one manuscript in preparation.
\end{itemize}%
}
\cventry
{2023 (6 months)}
{End-of-Studies Research Project}
{BioWooEB Research Group, CIRAD}
{Montpellier, France}
{}
{%
\emph{Optimization of a method for extracting proteins from dry rubber samples and development of a method for quantifying proteins using IR-TF}
\begin{itemize}[leftmargin=*,nosep]
    \item Performed SDS-PAGE analyses.
    \item Developed an infrared spectroscopy method for the characterisation
    of rubber proteins.
\end{itemize}%
}
\cventry
{2022--2023 (6 months)}
{Exchange Semester}
{Doshisha University}
{Kyoto, Japan}
{}
{%
\emph{Transdermal delivery of curcumin using an oil-in-water microemulsion based on isopropyl myristate.}
\begin{itemize}[leftmargin=*,nosep]
    \item Formulated and characterised curcumin-loaded oil-in-water microemulsions.
    \item Evaluated \emph{in vitro} transdermal delivery using Franz diffusion cell experiments.
    \item Performed spectrophotometric analyses to quantify curcumin permeation and release.
\end{itemize}%
}
\cventry
{2022 (2 months)}
{Research Internship}
{ADS Research Group, University of St Andrews}
{St Andrews, Scotland, UK}
{}
{%
\emph{Effect of substrates on the selectivity of the kinetic resolution of tertiary pyrazolone alcohols using chiral isothiourea catalysis}
\begin{itemize}[leftmargin=*,nosep]
    \item Synthesised pyrazolone derivatives and performed kinetic resolution
    using isothiourea catalysis.
    \item Purified compounds using column chromatography.
    \item Performed HPLC, NMR, and infrared spectroscopy analyses for
    structural characterisation and purity assessment.
\end{itemize}%
}
\cventry
{2020--2023}
{Research Project}
{École Nationale Supérieure de Chimie de Lille (ENSCL)}
{Lille, France}
{}
{%
\begin{itemize}[leftmargin=*,nosep]
\item Synthesised flame-retardant compounds for polylactide, including salen-based compounds and benzoxazines.
\item Performed reactive extrusion to manufacture PLA plates.
\end{itemize}%
}
%----------------------------------------------------------------------------------
% PUBLICATIONS
%----------------------------------------------------------------------------------
\section{Publications}

\subsection{Peer-reviewed publications}

\cvitem{First-author publications}{%
\textbf{Robichon, C.}, Villeneuve, P., Bohuon, P., Baréa, B.,
Barouh, N., Courtois, F., Fine, F., Durand, E. (2025).
\textit{Unveiling synergistic antioxidant combinations for
$\alpha$-tocopherol in emulsions: A spectrophotometric--mathematical
approach.}
\textit{Current Research in Food Science}, 11, 101134.
\href{https://doi.org/10.1016/j.crfs.2025.101134}
{\texttt{doi:10.1016/j.crfs.2025.101134}}
}

\cvitem{}{%
\textbf{Robichon, C.}, Durand, E., Gamaethiralalage, J. G.,
Bohuon, P., Baréa, B., Barouh, N., Courtois, F., Fine, F.,
de Smet, L. C. P. M., Villeneuve, P. (2026).
\textit{Decoding Tocopherol--Polyphenol Interactions in Oil-in-Water
Emulsions through Combined WIM--CAT and CV Assays.}
\textit{Current Research in Food Science}, 12, 101344.
\href{https://doi.org/10.1016/j.crfs.2026.101344}
{\texttt{doi:10.1016/j.crfs.2026.101344}}
}

\cvitem{Co-author publications}{%
Durand, E., Troncho, T., Koli, H., \textbf{Robichon, C.},
Villeneuve, P. (2026).
\textit{Focus on the Role of Mixed Micelles and Lipid Droplets in the
Oxidative Stability of Oil-in-Water Emulsions Using Size Distribution
Taylor Dispersion Analysis.}
\textit{Current Research in Food Science}, 12, 101369.
\href{https://doi.org/10.1016/j.crfs.2026.101369}
{\texttt{doi:10.1016/j.crfs.2026.101369}}
}

\cvitem{}{%
Smith, A. D., Westwood, M. T., Farah, A. O., Wise, H. B.,
Sinfield, M., \textbf{Robichon, C.}, Prindl, M. I.,
Cordes, D. B., Cheong, P. H.-Y. (2024).
\textit{Isothiourea-Catalysed Acylative Kinetic Resolution of
Tertiary Pyrazolone Alcohols.}
\textit{Angewandte Chemie International Edition}, e202407983.
\href{https://doi.org/10.1002/anie.202407983}
{\texttt{doi:10.1002/anie.202407983}}
}

\cvitem{Manuscript in preparation}{%
\textbf{Robichon, C.}, Durand, E., Silvestre, C., Vermorel, C.,
Roche, H., Troncho, T., Simon, L., Repellin, M.,
Dhuique-Mayer, C., Fine, F., Villeneuve, P.
\textit{Liposomes and Niosomes of $\alpha$-Tocopherol: Impact on
Bioaccessibility and Oxidative Stability in Oil-in-Water Emulsions.}
}

\cvitem{Submitted}{%
Villeneuve, P., \textbf{Robichon, C.}, Dhuique-Mayer, C.,
Durand, E. (2026).
\textit{Encapsulation of Vitamin E and Its Impact on Stability,
Delivery and Activity.}
}

%----------------------------------------------------------------------------------
% SCIENTIFIC COMMUNICATIONS
%----------------------------------------------------------------------------------
\clearpage
\section{Scientific Communications}
\cvitem{Oral}{%
20th Euro Fed Lipid Congress and Expo, Leipzig, Germany, 2025.
\emph{Design of a Spectrophotometric Assay Coupled with Mathematical Modeling
to Identify Optimal Antioxidant Synergies with Tocopherols in Oil-in-Water
Emulsions}.
}
\cvitem{Oral}{%
AOCS Annual Meeting \& Expo, Portland, USA, 2025.
\emph{Development of a Spectrophotometric Assay with Mathematical Modeling
to Unveil Optimal Combinations for Antioxidant Synergy with Tocopherols in
Oil-in-Water Emulsions}.
}
\cvitem{Poster}{%
5th International Symposium on Lipid Oxidation and Antioxidants,
Bologna, Italy, 2024.
\emph{A Rapid and Tunable Screening Method for Exploring Synergistic
Antioxidant Interactions with $\alpha$-Tocopherol in Oil-in-Water Emulsions}.
}
%----------------------------------------------------------------------------------
% SOFTWARE & CODE
%----------------------------------------------------------------------------------
\section{Software \& Code}
\cvitem{Model comparison}{MATLAB framework for comparing oxidation kinetic models. \href{https://github.com/camillerobichon30/Antioxidant-Kinetics-Model-Comparison}{\texttt{Available on GitHub}}.}
\cvitem{Antioxidant modelling}{MATLAB scripts for modelling antioxidant efficiency. \href{https://github.com/camillerobichon30/Weibull-Modelling-for-Antioxidant-Interaction-Efficiency}{\texttt{Available on GitHub}}.}
%----------------------------------------------------------------------------------
% SKILLS
%----------------------------------------------------------------------------------
\section{Technical Skills}
\cvitem{Formulation \& Systems}{%
Oil-in-water emulsions; nanoemulsions; liposomes; niosomes;
microfluidisation; reactive extrusion
}
\cvitem{In Vitro Models}{%
Static gastrointestinal digestion; vitamin E bioaccessibility;
transdermal delivery; Franz diffusion cells
}
\cvitem{Analytical Techniques}{%
HPLC; GC--MS; NMR; UV--Vis spectroscopy; infrared spectroscopy; DLS;
Taylor dispersion analysis; SDS-PAGE; cyclic voltammetry
}
\cvitem{Modelling \& Data Analysis}{%
MATLAB; kinetic modelling; model comparison; GraphPad Prism; VBA;
Microsoft Office 365
}

%----------------------------------------------------------------------------------
% LANGUAGES
%----------------------------------------------------------------------------------
\section{Languages}
\cvitem{French}{Native}
\cvitem{English}{Advanced professional proficiency -- TOEIC: 855/990}
\cvitem{Spanish}{Intermediate}
\cvitem{Japanese}{Beginner -- currently learning}
%----------------------------------------------------------------------------------
% INTERESTS
%----------------------------------------------------------------------------------
\section{Interests}
\cvitem{}{%
Detective novels; cycling; walking and hiking; concerts; Japanese language
and culture.
}
\end{document}
